\section{Conclusion}
\label{sec:conclusion}

The convite minimum-bidder rule of Lei 8.666/93, designed to
protect procedural competition, generates a behavioral residue
that procurement oversight bodies can detect at virtually no
marginal cost: firms that participate in many tenders without
ever winning, a pattern rationally dominated under any standard
entry model with positive bidding cost. We operationalize this
residue into a screening rule and validate it against CADE's
adjudication portfolio. Detection generalizes: AUC $=0.94$
against the headline ground truth, mean AUC $=0.954$
(SD $0.034$) across ten BEC product sectors, $2.6\times$
within-firm enrichment over the always-loser baseline. The
price-gap claim is weaker---a 3.6--7.7\% conditional
association concentrated where the regulatory cover-bidding
incentive operates---and the institutional design identifies
the 7.6\% premium under voluntary cover bidding within convite,
absent or negative under the binding constraint. Eight of the
twelve cartel cases were adjudicated after the sample window
closes, making the bulk of the validation prospective: the
screen anticipated firms that subsequent enforcement confirmed
years later.

We are explicit about what the design does not deliver. We do
not identify cover bidders firm-by-firm: conditional on
participation count, FL firms are statistically
indistinguishable from non-FL always-losers in CADE proximity
($p = 0.93$). We do not estimate causal overcharges: pre-event
coefficients in the FL-entry event study at $t = -2,-3$
preclude a fully causal reading. We do not measure welfare
losses beyond the illustrative $0.3$--$0.9\%$-of-spending range
in \ref{app:extensions} (Online), which assumes causality the
design does not establish. The Limitations section
(\S\ref{sec:limitations}) catalogs the residual threats and the
sensitivity bounds that constrain them.

\paragraph{Implementation pathway.}
The screen's contribution is operational, not methodological:
it is designed to run on the data oversight bodies already
maintain. We propose a three-stage deployment, ordered by
marginal data cost.

\begin{enumerate}
\item \textbf{Screen}. Apply the FL rule to contract-award
records: count tenders per firm in a defined window, flag
always-losers with participation above the
median $+\,1.5 \times$ IQR threshold. The computation runs in
seconds at the scale of millions of contracts and requires no
data transformation beyond winner/participant tallies. For
TCE-SP and analogous state audit bodies, this is a single
SQL query against the procurement warehouse they already query.

\item \textbf{Triage}. Re-score FL-flagged \emph{markets}
(PBU $\times$ item-group $\times$ year cells, not individual
firms: the $10\%$ firm-level persistence across our sample
halves means the flag identifies procurement environments,
not stable operatives) by winner-HHI and repeated co-bidding
pair counts. Prioritize markets with diversified winner
patterns and concentrated repeat partnerships---the
competitive-market FL subgroup that drives the price effect
in our network-split heterogeneity analysis
(Table~\ref{tab:network_split}, \S\ref{sec:results_network}).

\item \textbf{Investigate}. On the triaged subset, deploy
bid-level forensic tools---Bajari--Ye exchangeability and
conditional-independence tests, Imhof-style dispersion screens
\citep{imhof2019detecting, huber2019machine,
wallimann2023machine}, document discovery, and (where
available) leniency-driven testimony. Bid-level analysis is
expensive; the screen narrows the caseload to the markets where
the bid-level tools are most likely to pay off.
\end{enumerate}

The pathway is portable. For TCE-SP, CGE, and analogous
state-audit bodies, stages 1--2 require no new data
infrastructure. For CADE, the pathway offers an automated
triage layer for tip-driven investigations in settings where
bid-level microdata are not yet exported to analytical
warehouses. For OECD oversight bodies operating under
comparable minimum-bidder provisions
(EU Directive 2014/24/EU, US FAR 13.106-2; see
\S\ref{sec:institutional}), the same participation-only rule
applies. The screen's data parsimony makes it deployable
wherever a procurement system records participants and
winners, regardless of bid-level archiving.

\paragraph{Open questions.}
Three questions structure the agenda the paper opens.
\emph{Causality of the price effect}: the institutional design
identifies the regulatory channel under stated assumptions
(\S\ref{sec:empirical}), but a causal estimate of the
overcharge requires variation the within-window BEC sample
does not provide. The 2021 Brazilian reform (Lei 14.133/2021)
restricts convite and tightens habilitation; if post-2021 BEC
data become available, a difference-in-differences exploiting
the modality transition would test whether reducing the
regulatory cover-bidding incentive moves prices as the FL
diagnosis predicts. \emph{External validity}: replication in
ComprasNet (Brazilian federal procurement), other state-level
TCE platforms (TCE-MG, TCE-RS), and OECD platforms (Italy MEPA,
Mexico CompraNet) would establish whether AUC stability holds
out-of-sample. \emph{Sector--cartel correspondence}: the
inverse CNAE mapping in \S\ref{sec:cross_sector} recovered
eleven CNAE divisions for seven cartels but did not establish
a tight sector-to-cartel match; the BEC item-code dictionary
would close this and is the single most informative additional
data input.

Whether frequent losers are unlucky competitors or cartel
operatives, their presence marks procurement environments that
warrant a closer look---and the closer look can begin with
data oversight bodies already have.
